\section{Reflexión final}
\label{sec:concl_reflexion}
% Ficha 7.5. Única sección en primera persona (comentario del tutor del 13/11; modelo: tesina de referencia, 5.4).

Desde un punto de vista personal, este trabajo final de carrera fue una experiencia de aprendizaje muy enriquecedora.
Me permitió integrar los conocimientos adquiridos a lo largo de la carrera con la investigación en un área que avanza
muy rápido, como la de los modelos visión--lenguaje, y llevarlos a un problema clínico concreto. El trabajo comenzó con
la hipótesis de que un único modelo podría diagnosticar la lesión y redactar el reporte, y los resultados me obligaron
a replantearla. Esa experiencia me enseñó que un resultado negativo, bien medido, también orienta el diseño de un
sistema.

También aprendí que, en un problema clínico, las decisiones de evaluación pesan tanto como la elección del modelo: la
forma de dividir los datos, la métrica o el criterio con que se cuenta un acierto pueden cambiar las conclusiones. Por
último, considero que la inteligencia artificial puede aportar a la práctica dermatológica en la medida en que se
conciba como apoyo al especialista y se evalúe con el mismo rigor que cualquier otra herramienta clínica, y este
trabajo fue una oportunidad para recorrer ese camino.
